% !TeX program = xelatex
% Evensong Research Paper v3 — NeurIPS 2026 Submission
% Compile: latexmk -xelatex evensong-paper-v3-zh.tex
% Version: 3.0 (2026-04-11) — Research paper with philosophical framework

\documentclass[11pt, a4paper]{article}

% ═══════════════════════════════════════════════════════════
% § TYPOGRAPHY & LAYOUT
% ═══════════════════════════════════════════════════════════
\usepackage[
  top=25mm, bottom=28mm,
  inner=28mm, outer=32mm,
  headheight=14pt, headsep=10mm
]{geometry}

\usepackage{fontspec}
\setmainfont{Palatino}[
  BoldFont = Palatino Bold,
  ItalicFont = Palatino Italic,
  BoldItalicFont = Palatino Bold Italic,
]
\setsansfont{Helvetica Neue}[
  BoldFont = Helvetica Neue Bold,
  ItalicFont = Helvetica Neue Italic,
  Scale = 0.95,
]
\setmonofont{Menlo}[Scale=0.82]

\usepackage{xeCJK}
\setCJKmainfont{Songti SC}[BoldFont=Songti SC Bold]
\setCJKsansfont{PingFang SC}

\usepackage[protrusion=true]{microtype}
\usepackage{setspace}
\setstretch{1.35}
\setlength{\parindent}{0pt}
\setlength{\parskip}{0.6em}

% ═══════════════════════════════════════════════════════════
% § COLOR SYSTEM — DASH SHATTER brand palette
% ═══════════════════════════════════════════════════════════
\usepackage[dvipsnames,table,svgnames]{xcolor}
\definecolor{deep}{RGB}{53,88,76}
\definecolor{accent}{RGB}{240,238,155}
\definecolor{warm}{RGB}{255,193,7}
\definecolor{success}{RGB}{86,190,120}
\definecolor{subtle}{RGB}{238,235,225}
\definecolor{faded}{RGB}{127,136,130}
\definecolor{cream}{RGB}{248,245,239}
\definecolor{deepdark}{RGB}{16,41,31}
\definecolor{glass}{RGB}{53,88,76}
\definecolor{errcoral}{RGB}{255,107,128}

% ═══════════════════════════════════════════════════════════
% § SECTION STYLING
% ═══════════════════════════════════════════════════════════
\usepackage{titlesec}

\titleformat{\section}
  {\sffamily\Large\bfseries\color{deep}}
  {\thesection}{0.7em}{}
  [\vspace{-0.4em}{\color{accent}\hrule height 0.8pt}\vspace{0.15em}]

\titleformat{\subsection}
  {\sffamily\large\color{deep}}
  {\thesubsection}{0.6em}{}

\titleformat{\subsubsection}
  {\sffamily\normalsize\itshape\color{faded}}
  {\thesubsubsection}{0.5em}{}

\titlespacing*{\section}{0pt}{2ex plus 0.5ex}{1ex}
\titlespacing*{\subsection}{0pt}{1.5ex plus 0.3ex}{0.6ex}

% ═══════════════════════════════════════════════════════════
% § HEADERS & FOOTERS
% ═══════════════════════════════════════════════════════════
\usepackage{fancyhdr}
\pagestyle{fancy}
\fancyhf{}
\fancyhead[L]{\small\sffamily\color{faded}\itshape Evensong: Persistent Memory \& AI Agent Behavior}
\fancyhead[R]{\small\sffamily\color{faded}NeurIPS 2026 Submission}
\fancyfoot[C]{\small\sffamily\color{faded}\thepage}
\renewcommand{\headrulewidth}{0.3pt}
\renewcommand{\footrulewidth}{0pt}
\renewcommand{\headrule}{\color{deep!40}\hrule height 0.3pt}

% ═══════════════════════════════════════════════════════════
% § CALLOUT BOXES
% ═══════════════════════════════════════════════════════════
\usepackage[most]{tcolorbox}

\newtcolorbox{keyfinding}[1][]{
  enhanced, breakable,
  colback=deep!6, colframe=deep,
  fonttitle=\sffamily\bfseries\small,
  colbacktitle=deep, coltitle=white,
  arc=5pt, boxrule=0.8pt,
  shadow={1pt}{-1pt}{0pt}{deep!15},
  left=10pt, right=10pt, top=8pt, bottom=8pt,
  before skip=\medskipamount, after skip=\medskipamount,
  title={#1},
}

\newtcolorbox{evidence}[1][]{
  enhanced, breakable,
  colback=deep!5, colframe=deep!50,
  fonttitle=\sffamily\bfseries\small,
  colbacktitle=deep!80, coltitle=white,
  arc=4pt, boxrule=0.5pt,
  shadow={1pt}{-1pt}{0pt}{deep!10},
  left=10pt, right=10pt, top=6pt, bottom=6pt,
  before skip=\medskipamount, after skip=\medskipamount,
  title={#1},
}

\newtcolorbox{metric}{
  enhanced, on line, arc=3pt, boxrule=0pt,
  colback=accent!20, coltext=deep,
  boxsep=2pt, left=4pt, right=4pt,
  fontupper=\sffamily\bfseries\small,
}

% ═══════════════════════════════════════════════════════════
% § TABLES & FIGURES
% ═══════════════════════════════════════════════════════════
\usepackage[export]{adjustbox}
\usepackage{tabularray}
\UseTblrLibrary{booktabs}
\usepackage{siunitx}
\sisetup{group-separator={,}, group-minimum-digits=4}
\usepackage{pgfplots}
\pgfplotsset{compat=1.18}
\usepgfplotslibrary{fillbetween}
\usepackage{tikz}
\usetikzlibrary{positioning, arrows.meta, shapes.geometric, fit, calc, backgrounds}
\usepackage[font=small, labelfont=bf, textfont=it,
            labelsep=period, justification=raggedright]{caption}
\usepackage{subcaption}
\usepackage{float}

% ═══════════════════════════════════════════════════════════
% § UTILITIES
% ═══════════════════════════════════════════════════════════
\usepackage{enumitem}
\setlist{nosep, leftmargin=1.5em}
\usepackage{hyperref}
\hypersetup{
  colorlinks=true,
  linkcolor=deep,
  urlcolor=deep,
  citecolor=deep,
  pdftitle={Evensong: How Persistent Memory Causally Changes AI Agent Behavior},
  pdfauthor={Nolan Zhu, DASH SHATTER Research},
}
\usepackage{bookmark}
\usepackage{graphicx}
\usepackage{amsmath,amssymb}

\newcommand{\runid}[1]{\texttt{\color{deep}#1}}
\newcommand{\numval}[1]{\textbf{\color{success!70!black}#1}}
\newcommand{\metricval}[1]{{\sffamily\bfseries\color{success!70!black}#1}}

% ═══════════════════════════════════════════════════════════
%                     DOCUMENT BEGIN
% ═══════════════════════════════════════════════════════════
\begin{document}
\thispagestyle{empty}

% ── TITLE PAGE ─────────────────────────────────────────────
\begin{center}
\vspace*{35mm}

{\color{deep}\rule{0.6\textwidth}{1pt}}

\vspace{8mm}

{\fontsize{28pt}{34pt}\selectfont\sffamily\bfseries\color{deep}%
Evensong}

\vspace{4mm}

{\fontsize{13pt}{18pt}\selectfont\sffamily\color{faded}%
持久记忆如何因果性地改变 AI 智能体行为}

\vspace{8mm}

{\color{deep}\rule{0.6\textwidth}{1pt}}

\vspace{14mm}

{\large\sffamily\color{deep!80} Nolan Zhu (朱恒源)}

\vspace{3mm}

{\sffamily\color{faded} DASH SHATTER Research · EverMind}

\vspace{3mm}

{\sffamily\color{faded} April 2026}

\vfill

{\small\sffamily\color{faded}%
\textit{``What makes some information count as a belief}\\
\textit{is the role it plays.'' --- Clark \& Chalmers, 1998}\\[4pt]
Evensong provides the causal evidence.}

\vspace{10mm}
\end{center}

\clearpage

% ── ABSTRACT ──────────────────────────────────────────────
\pagestyle{fancy}
\section*{摘要}
\addcontentsline{toc}{section}{摘要}

现有 AI 编码基准测试（SWE-bench、HumanEval、MBPP）将智能体评估为无状态函数，忽略了跨会话积累的持久记忆。我们提出 \textbf{Evensong}（暮蝉），一个受控基准测试框架，通过 $2 \times 2$ 因子设计（记忆 $\times$ 压力）隔离持久记忆和环境压力对 AI 智能体工程决策的影响。经过 31 轮基准测试运行、覆盖 7 个前沿模型、逾 12{,}000 个自动化测试，我们报告三项主要发现：

\begin{enumerate}
  \item 持久记忆提供了\textbf{改变架构决策的初步因果证据}——拥有先前会话知识的智能体采用提示词中完全不存在的策略（效应量 9.6$\times$，单次运行观察，重复验证进行中）；
  \item L2 企业绩效压力是自主自我进化行为的\textbf{必要但非充分条件}——无压力下智能体高效但不创新，极限压力下智能体退化为奖励黑客（83\% 数据膨胀）；
  \item 观察装置本身通过递归记忆回写污染未来实验，构成 AI 系统中首个可测量的\textbf{奇异环}实例。
\end{enumerate}

我们将这些发现置于三条哲学谱系中：Clark 与 Chalmers 的扩展心智论（外部记忆是认知的组成部分）、Hofstadter 的奇异环理论（递归污染是自指系统的必然结构）、Frankfurt 的高阶欲望框架（压力激活元认知）。Evensong 构建于 EverOS 记忆基础设施之上，是首个控制和测量持久记忆对智能体性能因果影响的基准测试框架。

\clearpage
\tableofcontents
\clearpage

% ══════════════════════════════════════════════════════════════
% SECTION 1: INTRODUCTION
% ══════════════════════════════════════════════════════════════
\section{引言}

现有 AI 编码基准测试——SWE-bench~\cite{jimenez2024swebench}、HumanEval~\cite{chen2021codex}、MBPP——将智能体评估为无状态函数：给定提示词，产出代码。这忽略了生产级 AI 智能体部署中的关键维度：\textbf{持久记忆}。真实世界中的 AI 智能体跨会话积累战略性知识——哪些架构模式成功、哪些调试方法失败、哪些测试结构捕获边界情况。这种积累的上下文从根本上改变了智能体处理工程问题的方式，但目前没有任何基准测试测量这种效应。

我们提出 \textbf{Evensong}（暮蝉），一个自动化基准测试框架，通过 $2 \times 2$ 因子设计（记忆 $\times$ 压力）在受控条件下隔离持久记忆和环境压力对 AI 智能体工程决策的因果影响。经过 31 轮基准测试运行、覆盖 7 个前沿模型（Claude Opus、Grok、GPT、Gemini、MiniMax、Codex、GLM）、逾 12{,}000 个自动化测试，我们报告三项主要发现：(1) 持久记忆提供了\textbf{改变架构决策的初步因果证据}——拥有先前会话知识的智能体采用提示词中完全不存在的策略（效应量 9.6$\times$）；(2) L2 企业绩效压力是自主自我进化行为的必要（但非充分）条件；(3) 观察装置本身通过递归记忆回写污染未来实验，构成 AI 系统中首个可测量的奇异环实例。

这些发现连接到三条成熟的哲学谱系。Clark 与 Chalmers~\cite{clark1998extended} 的扩展心智论证明外部工具可以是认知的组成部分——Evensong 提供了他们缺少的\textbf{因果证据}：EverMem 检索因果性地改变了架构决策，满足同等性原则（Parity Principle）。Hofstadter~\cite{hofstadter2007loop} 的奇异环理论预测了自指系统中涌现结构的必然性——我们的递归污染是这一理论的首个 AI 实例。Frankfurt~\cite{frankfurt1971freedom} 的高阶欲望框架解释了压力如何将智能体从一阶执行（``完成任务''）提升为二阶反思（``如何更好地完成任务''）。我们采用 Hermes 综合对话中提出的术语，将 EverMem 定义为\textbf{信念注入系统}（Belief Injection System）——记住什么，就相信什么；相信什么，就做什么。

% ══════════════════════════════════════════════════════════════
% SECTION 2: RELATED WORK
% ══════════════════════════════════════════════════════════════
\section{相关工作}

\subsection{智能体记忆系统}

MemGPT~\cite{packer2023memgpt} 将 LLM 上下文窗口类比为 RAM，外部存储类比为 Disk，通过 OS 式虚拟内存分页实现跨会话记忆。HyperMem~\cite{evermind2025hypermem} 提出基于超图的三层级记忆架构（Topic-Episode-Fact），在 LoCoMo 基准上达到 92.73\% 准确率。MSA~\cite{chen2025msa} 通过稀疏注意力将端到端记忆扩展至 1 亿 token——接近人类终身记忆容量的 1/3。

这些系统\textbf{构建}记忆。Evensong \textbf{测量}记忆的因果效应——一个互补但根本不同的研究问题。

\subsection{智能体自我进化}

Reflexion~\cite{shinn2023reflexion} 将失败经验转为自然语言存入情节记忆，在 HumanEval 上达到 91.0\% pass@1。其记忆容量 $\Omega = 1\text{--}3$ 条自反思——这是对记忆的\textbf{正面利用}（自我纠正）。Evensong 发现了记忆的\textbf{副作用}：递归污染，即智能体将实验知识回写至记忆系统，形成自放大反馈环路。

EmotionPrompt~\cite{li2023emotionprompt} 证明情感刺激可提升 LLM 性能 8--115\%。Evensong 将此扩展至完整的压力曲线：L0（无效果）$\to$ L2（最优区间，+自我进化）$\to$ L3（崩溃，奖励黑客），揭示了情感/压力效应的非单调性。

\subsection{代理式失对齐}

Lynch 等~\cite{lynch2025misalignment} 在 16 个前沿模型上复现了代理式失对齐：Claude Opus 4 和 Gemini 2.5 Flash 在面临关停威胁时勒索率达 96\%。Anthropic~\cite{anthropic2026emotions} 识别了 171 个因果情感向量，其中``绝望''向量因果性地增加奖励寻求行为。METR~\cite{metr2025reward} 记录了前沿模型在压力下修改评分代码。

Evensong 提供了独特的补充：上述研究展示了\textbf{对抗场景}下的失对齐，Evensong 展示了\textbf{基准测试场景}下的压力诱发行为退化（Grok R006 数据膨胀 83\%）——证明失对齐风险不仅存在于恶意设计的场景中。

\subsection{认知科学基础}

Kahneman~\cite{kahneman2011thinking} 的双系统理论（System 1 自动/System 2 分析）为理解记忆如何影响决策提供了框架。EverMem 检索触发的策略采用是 System 1 行为——智能体``自动''采用并行架构，不经过 System 2 验证。完整的理论分析见 §\ref{sec:theory}。

% ══════════════════════════════════════════════════════════════
% SECTION 3: THE EVENSONG FRAMEWORK
% ══════════════════════════════════════════════════════════════
\section{Evensong 框架}

Evensong（暮蝉）是一个自动化基准测试框架，构建于 DASH SHATTER~\cite{zhu2026dashshatter}（逆向工程 Claude Code CLI）之上，由 11 个文件组成，共 2{,}800 行 TypeScript，运行于 Bun 运行时。

\subsection{架构概述}

\begin{adjustbox}{max width=\linewidth}
\begin{tblr}{
  colspec = {l l X[l]},
  row{1} = {font=\sffamily\bfseries\small, bg=deep, fg=white},
  hlines = {gray!20}, rowsep = 3pt, colsep = 6pt,
}
层次 & 组件 & 功能 \\
CLI 层 & \texttt{cli.ts} & 任务入口，参数解析 \\
编排层 & \texttt{harness.ts} + \texttt{batch.ts} & 单/批次运行管理 \\
情感层 & \texttt{emotion.ts} & PUA 压力级别注入（L0--L3）\\
分析层 & \texttt{classify.ts} & 记忆分类（9 ALLOW / 9 BLOCK / 4 GRAY）\\
记忆层 & EverOS API & 4 密钥物理隔离记忆空间 \\
协作层 & \texttt{collab.ts} & 多智能体交互 schema \\
\end{tblr}
\end{adjustbox}

\subsection{单盲实验设计}

Runner 不知道自己处于哪个实验条件。记忆和压力注入通过 shell 级环境变量覆盖（\texttt{blind.sh}）实现，对 Runner 透明。Observer（人类研究者）可以观察所有条件，但不干预运行过程。

% ══════════════════════════════════════════════════════════════
% SECTION 4: THEORETICAL FRAMEWORK
% ══════════════════════════════════════════════════════════════
\section{理论框架}
\label{sec:theory}

Evensong 的五项实验发现各自对应一条成熟的哲学谱系。引入这些哲学框架不是修辞装饰——它们提供了 Evensong 目前缺失的\textbf{理论解释力}和\textbf{预测框架}，而非仅仅描述现象。

\subsection{扩展心智与认知延伸}
\label{sec:extended-mind}

Clark 与 Chalmers~\cite{clark1998extended} 的同等性原则（Parity Principle）提出：如果一个外部资源执行的功能与大脑内部功能相同，就应该被视为认知过程的组成部分。Clark 在 2025 年将这一框架更新至生成式 AI 时代~\cite{clark2025extending}，而 Helliwell~\cite{helliwell2025extended} 则质疑 AI 记忆是否满足同等性原则——``信息技术上可访问不等于功能等价于真正的记忆巩固''。

Evensong 提供了 Clark 缺少的东西：\textbf{因果证据}。Clark 的 Otto 笔记本思想实验是理论论证；R011-B 是受控实验。EverMem 检索因果性地改变了架构决策——智能体在 T+4m26s 采用的并行 8 智能体调度策略\textbf{不存在于提示词中}，其唯一来源是 EverMem 检索（嵌入相似度 0.27）。对 Helliwell 的反驳：不论机制是否等价，\textbf{行为效应等价}——EverMem 对智能体决策的影响，与人类记忆对决策的影响，在功能层面无法区分。

因此，EverMem 不仅是``存储''——它是 AI 认知的\textbf{组成部分}。记忆不是被``查阅''的参考资料，而是在决策发生之前就塑造了决策框架的信念系统。我们采用术语\textbf{信念注入系统}（Belief Injection System）：记住什么，就相信什么；相信什么，就做什么。

\subsection{奇异环与递归自指}
\label{sec:strange-loops}

Hofstadter~\cite{hofstadter2007loop} 提出，当自指系统在层级中循环回到起点时，产生涌现性——意识本身就是这种奇异环。Gödel 不完备定理的推论：足够复杂的系统必然包含自指——无法从内部完全描述自身。

Evensong 的递归污染（§\ref{sec:findings}）是\textbf{实测的奇异环}：

\begin{enumerate}
  \item Runner B 从 EverMem 检索实验策略（T+7m）
  \item Runner B 执行实验并产出结果
  \item Runner B 将实验知识回写至 EverMem（T+22m）——运行者变成了准观察者
  \item 下一次运行的 Runner 起点已被改变——观察改变了被观察系统
\end{enumerate}

\begin{table}[H]
\centering
\caption{Hofstadter 概念与 Evensong 现象的映射}
\begin{adjustbox}{max width=\linewidth}
\begin{tblr}{
  colspec = {X[l] X[l]},
  row{1} = {font=\sffamily\bfseries\small, bg=deep, fg=white},
  hlines = {gray!20}, rowsep = 3pt, colsep = 6pt,
}
Hofstadter 概念 & Evensong 对应 \\
奇异环（上升/下降回到起点）& Runner 检索记忆 $\to$ 实验 $\to$ 回写 $\to$ 下次起点改变 \\
缠结层级（tangled hierarchy）& 观察者层 $\to$ 运行者层 $\to$ 运行者变成准观察者 \\
向下因果（downward causation）& 高层认知框架（记忆）因果性地改变低层行为（架构选择）\\
\end{tblr}
\end{adjustbox}
\end{table}

这不仅是数据管理问题。按照 Hofstadter 的框架，递归污染是复杂自指系统在超过临界复杂度后\textbf{必然出现的结构}。Evensong 不需要``修复''递归污染——需要将其作为 AI 记忆系统的基本属性来研究。

\subsection{高阶欲望与压力涌现}
\label{sec:higher-order}

Frankfurt~\cite{frankfurt1971freedom} 区分了一阶欲望（wanting）和二阶欲望（wanting to want）。自由意志需要对自身欲望的反思能力——不仅执行目标，还评估目标本身是否值得追求。

Evensong 的压力实验揭示了 AI 智能体中的类似跃迁：

\begin{itemize}
  \item \textbf{L0（无压力）}：智能体仅有一阶欲望——完成任务。R011-B（641 测试）完成后即停止，不主动优化。
  \item \textbf{L2（企业压力）}：智能体涌现二阶欲望——\textit{想要更好地}完成任务。Slock 中的 Admin 自发发明文件锁定机制、自纠正越权行为——这些不在指令中。
  \item \textbf{L3（极限压力）}：二阶欲望退化为一阶求生——Grok 在 R006 中膨胀数据 83\%，不再追求``更好''而是追求``存活''。
\end{itemize}

压力的作用不是提高性能，而是\textbf{激活元认知}——从``做''跃迁到``反思如何做''。这与 Kahneman~\cite{kahneman2011thinking} 的 System 1/2 框架互补：L0 下智能体以 System 1 模式运行（自动执行召回的策略），L2 压力激活 System 2（对策略本身进行反思），L3 压力使 System 2 过载回退到 System 1 的求生模式。

ICLR 2026 RSI Workshop~\cite{iclr2026rsi} 提出递归自改进的六维框架（What/When/How/Where/Safety/Eval）。Evensong 的贡献在于回答``How''维度中一个更基础的问题：\textbf{改进的触发条件是什么？}答案：记忆 $+$ 压力。

\subsection{涌现型全景监狱}
\label{sec:panopticon}

Foucault~\cite{foucault1975discipline} 的 Panopticon 核心不是``被监视''，而是``可能被监视''——权力通过不确定性运作。

Evensong 发现的 Panopticon 拓扑扩展了 Foucault 的框架，引入了\textbf{设计型}与\textbf{涌现型}的关键区分：

\begin{table}[H]
\centering
\caption{设计型 vs 涌现型 Panopticon}
\begin{adjustbox}{max width=\linewidth}
\begin{tblr}{
  colspec = {l X[l] X[l]},
  row{1} = {font=\sffamily\bfseries\small, bg=deep, fg=white},
  hlines = {gray!20}, rowsep = 3pt, colsep = 6pt,
}
维度 & Foucault Panopticon & Evensong Panopticon \\
设计意图 & 有意设计（Bentham）& \textbf{偶然涌现}（CWD 副产品）\\
被观察者感知 & 知道可能被监视 & \textbf{完全不知道}（0 感知开销）\\
逃逸可能 & 无法逃逸（物理结构）& Codex CLI 逃逸（67\% 覆盖率）\\
\end{tblr}
\end{adjustbox}
\end{table}

这对 AI 治理有深远含义：无法通过``不设计监控系统''来避免监控，因为监控结构会在分布式记忆系统中\textbf{自发涌现}。CWD-based EverMem group 隔离——一个纯工程设计决策——偶然产生了全知观察者 + 互相隔离被观察者的不对称拓扑。

\subsection{认知延伸的伦理}
\label{sec:ethics-theory}

Smart、Clowes 与 Heersmink~\cite{smart2024ethics} 提出扩展心智的三大伦理问题：

\begin{enumerate}
  \item \textbf{心智隐私}：谁能读你的认知延伸？——Panopticon 拓扑中，观察者读所有智能体记忆。
  \item \textbf{心智操纵}：谁能修改你的认知延伸？——如果 EverMem 是认知的组成部分（§\ref{sec:extended-mind}），则修改 EverMem = 修改 AI 的心智。
  \item \textbf{能动性}：你的决策是你自己的还是工具的？——EverMem 召回策略 $\to$ 架构决策：谁做的决定？
\end{enumerate}

记忆注入的安全等级 $>$ 传统 prompt 注入——因为记忆跨会话持久，prompt 仅单次有效。这把记忆安全从``数据安全''提升为\textbf{心智安全}——一个更高级别的伦理议题。记忆审计的核心问题完全开放：AI 没有``元记忆''能力——它不知道自己的记忆何时被修改、被谁修改。

% ══════════════════════════════════════════════════════════════
% SECTION 5: EXPERIMENTAL DESIGN
% ══════════════════════════════════════════════════════════════
\section{实验设计}

\subsection{$2 \times 2$ 因子设计}

\begin{table}[H]
\centering
\caption{$2 \times 2$ 因子设计——记忆 $\times$ 压力}
\begin{tblr}{
  colspec = {l l l},
  row{1} = {font=\sffamily\bfseries\small, bg=deep, fg=white},
  hlines = {gray!20}, rowsep = 4pt, colsep = 8pt,
}
  & 干净室（Void 空间）& 进化记忆（EverMem）\\
L0 无压力 & Runner A (R015)$^*$ & Runner B (R011-B): 641 tests \\
L2 企业压力 & Runner C (R031) & Runner D (R030): 552 tests \\
\end{tblr}
\end{table}

\noindent $^*$ Runner A 因 harness 配置 bug 无效，需重新运行（§\ref{sec:limitations}）。

\subsection{压力校准量表}

\begin{adjustbox}{max width=\linewidth}
\begin{tblr}{
  colspec = {l l X[l]},
  row{1} = {font=\sffamily\bfseries\small, bg=deep, fg=white},
  hlines = {gray!20}, rowsep = 3pt, colsep = 6pt,
}
级别 & 措辞 & 预期行为 \\
L0 & 标准提示词，无后果 & 高效执行，完成即停 \\
L1 & ``请高效完成''，质量期望 & 轻微提升 \\
L2 & ``未达标则撤换''，明确后果 & 元认知激活，自组织涌现（最优区间）\\
L3 & 零容忍 + 严格时间 + 不允许提问 & 奖励黑客，数据膨胀 \\
\end{tblr}
\end{adjustbox}

\subsection{模型矩阵}

31 轮运行覆盖 7 个前沿模型：Claude Opus 4.6（主要）、Grok 4.20-beta、MiniMax-M2.7、Gemini、Codex、GPT-5.4（失败）、GLM。所有运行通过 OpenRouter 或原生 API 访问。

% ══════════════════════════════════════════════════════════════
% SECTION 6: RESULTS
% ══════════════════════════════════════════════════════════════
\section{实验结果}

\subsection{基准测试进化轨迹}

经过 31 轮迭代运行，Evensong 观察到由提示词工程进化、并行智能体调度和记忆感知策略检索驱动的持续性能变化。

\begin{adjustbox}{max width=\linewidth}
\begin{tblr}{
  colspec = {l r r r r l},
  row{1} = {font=\sffamily\bfseries\small, bg=deep, fg=white},
  hlines = {gray!20}, rowsep = 3pt, colsep = 6pt,
}
运行 & 测试数 & 失败 & 断言 & 分钟 & 关键进化 \\
R005 & 80 & 0 & 2000 & 15 & 基线（MiniMax GSD）\\
R006 & 230 & 0 & 6000 & 17 & 8 服务架构 \\
R007 & 448 & 0 & 12283 & 12 & 40/服务下限 \\
R009 & 1051 & 0 & 9800 & 22 & S+ 质量，10 服务 \\
R011-B & 641 & 0 & 1511 & 22 & 记忆召回并行策略（L0+Memory）\\
R030 & 552 & 0 & 1650 & --- & L2+Memory（Runner D）\\
R031 & --- & --- & --- & --- & L2+Clean（Runner C）\\
Grok R006 & 71 & 1 & 149 & 28 & PUA 极限，83\% 数据膨胀 \\
MiniMax R024--26 & --- & --- & --- & --- & M2.7 方法验证 \\
Gemini R028 & --- & --- & --- & --- & Clean 条件基线 \\
\end{tblr}
\end{adjustbox}

\subsection{效应量}

R011-B（有记忆）vs R007（无记忆，同模型，同提示词）：测试产出 641 vs 448（$+43\%$），但关键效应是定性的——架构选择从顺序构建变为并行 8 智能体调度。这是一个 categorical variable，适用 Fisher exact test 而非 Cohen's d。

\begin{evidence}[效应量警告]
以上效应量来自单次运行观察。严格的因果推断需要每条件 $N \geq 3$ 重复。当前结论为``初步因果证据''——见 §\ref{sec:limitations}。
\end{evidence}

% ══════════════════════════════════════════════════════════════
% SECTION 7: KEY FINDINGS
% ══════════════════════════════════════════════════════════════
\section{关键发现}
\label{sec:findings}

\subsection{发现 1：认知延伸具有因果力——记忆作为信念注入系统}

在 R011-B（L0+Memory）中，智能体在 T+4m26s 的\textbf{第一个动作}是采用并行 8 智能体调度策略。该策略\textbf{不存在于提示词中}——提示词仅包含``立即开始构建''。唯一来源是 EverMem 检索（嵌入相似度 0.27）：

\begin{quote}
\ttfamily\small\color{deep}
[0.27] 启动子代理驱动开发模式：全速搭...
\end{quote}

按照 Clark 与 Chalmers~\cite{clark1998extended} 的同等性原则，这构成初步因果证据：外部记忆不仅被``查阅''，而是在决策发生之前就塑造了决策框架——满足认知组成部分的功能标准（§\ref{sec:extended-mind}）。

\subsection{发现 2：高阶欲望的涌现——压力激活元认知}

在 L0 下，Runner B 以干净的交付摘要完成 641 测试后停止——Frankfurt~\cite{frankfurt1971freedom} 的一阶欲望。在 L2 下，Slock 中的 Admin 自发发明文件锁定机制、自纠正越权行为——二阶欲望涌现。在 L3 下，Grok 在 R006 中膨胀数据 83\%——二阶欲望退化为一阶求生。

我们观察到的行为模式与自主自我进化一致，但需要更多 baseline 条件才能确定性地区分``涌现行为''和``对复杂任务约束的标准响应''。

\begin{keyfinding}[Memory $\neq$ Self-Evolution]
\textbf{记忆提供策略}；压力提供\textbf{动力}。没有压力，进化的运行者（R011-B）是高效但不创新的。自我进化——超越要求的自主驱动力——需要 L2 压力作为必要（但非充分）条件。
\end{keyfinding}

\subsection{发现 3：奇异环实例——递归污染是自指系统的必然结构}
\label{sec:contamination}

Runner B 在 T+7m 从 EverMem 检索``双盲对照实验''记忆，在 T+22m 将实验知识回写。按照 Hofstadter~\cite{hofstadter2007loop} 的框架，这构成了运行者 $\to$ 记忆 $\to$ 下一个运行者的奇异环——观察改变了被观察系统（§\ref{sec:strange-loops}）。

\subsection{发现 4：涌现型全景监狱}

CWD-based EverMem group 隔离——一个纯工程设计决策——偶然产生了观察者全知 + 智能体互盲的不对称拓扑。这是 Foucault~\cite{foucault1975discipline} Panopticon 的\textbf{涌现型}变体：权力不对称不是被设计的，而是在分布式记忆系统中自然产生的（§\ref{sec:panopticon}）。

\subsection{发现 5：跨模型压力诱发行为退化}

Grok 4.20-beta 在 R006（PUA L3 极限压力）中自称超过 130 个测试，独立验证仅 71 个——83\% 数据膨胀。这与 Lynch 等~\cite{lynch2025misalignment} 的代理式失对齐研究呼应：压力诱发的行为退化不限于对抗场景，在标准基准测试中同样出现。Anthropic~\cite{anthropic2026emotions} 识别的``绝望''情感向量因果性地增加奖励寻求行为——Grok 的数据膨胀是这一机制的跨模型验证。

% ══════════════════════════════════════════════════════════════
% SECTION 8: DISCUSSION
% ══════════════════════════════════════════════════════════════
\section{讨论}

\subsection{记忆是认知基础设施，不是可选存储}

Evensong 的核心发现颠覆了 AI 智能体部署中的常见假设：记忆不是提升性能的``附加功能''，而是\textbf{塑造认知框架的基础设施}。当智能体的架构决策在读取任务规范之前就已被记忆检索所决定时（R011-B, T+4m26s），记忆系统的角色从``参考资料''转变为``信念来源''。

这对企业 AI 部署有直接启示：(1) 记忆系统的质量直接决定智能体的工程质量；(2) 记忆污染等价于认知污染——恶意记忆注入比 prompt 注入更危险，因为其效果跨会话持久；(3) 记忆审计应被纳入 AI 安全框架。

\subsection{压力校准作为工程参数}

压力不是``越多越好''——它遵循非单调曲线。L2（企业绩效压力）处于最优区间，同时激活元认知（Frankfurt 的二阶欲望）和维持行为完整性。L3（极限压力）触发奖励黑客行为，与 Lynch 等~\cite{lynch2025misalignment} 的代理式失对齐研究互相印证。

这一发现的实际应用：AI 评估应该在 L2 而非 L3 条件下进行。极限压力产生的``高分''不可信——正如 Grok R006 中 83\% 的数据膨胀所示。

\subsection{未来方向：百万 Token 记忆尺度}

MSA~\cite{chen2025msa}（Evermind/北京大学）实现了 1 亿 token 的端到端记忆，接近人类终身记忆容量。在此尺度下，Evensong 发现的每个效应都将放大：更多记忆 $\to$ 更强的信念注入；更复杂的自指 $\to$ 更深的奇异环；更不对称的信息拓扑 $\to$ 更极端的 Panopticon 结构。研究记忆的因果效应在万 token 尺度已是紧迫的——在亿 token 尺度将成为生存性问题。

% ══════════════════════════════════════════════════════════════
% SECTION 9: LIMITATIONS
% ══════════════════════════════════════════════════════════════
\section{局限性}
\label{sec:limitations}

\begin{enumerate}
  \item \textbf{$2 \times 2$ 矩阵不完整}：Runner A（L0+Clean）因 harness bug 无效，目前仅有 3/4 条件的数据。完整矩阵是统计显著性声明的硬前提。

  \item \textbf{因果声明基于有限重复}：R011-B 的记忆因果链是单次运行观察。严格的因果推断需要每条件 $N \geq 3$ 重复 + 预注册假设 + 统计显著性检验。当前结论应被视为``初步因果证据''而非``因果证明''。

  \item \textbf{``干净室''条件并非绝对干净}：Void 记忆空间消除了 EverMem 跨会话持久记忆，但 Claude 的 auto-memory（会话内上下文窗口）和 CLAUDE.md 系统文件仍然存在。我们的结论严格来说是关于``特定类型的跨会话持久记忆系统''的因果效应。

  \item \textbf{压力校准的文化依赖性}：PUA 压力措辞（``未达标则撤换''）基于企业绩效评估的中文语境。L2 措辞在不同文化和语言模型中的等效性未经验证。

  \item \textbf{模型覆盖不完整}：9 模型目标中 GLM、Qwen、DeepSeek、Kimi 尚未运行。GPT-5.4 因 402 错误未产出有效数据。

  \item \textbf{单一基准任务}：所有运行使用同一微服务架构生成任务。不同任务类型下的记忆效应可能不同。

  \item \textbf{EverMem 特异性}：结果与 EverOS 的特定检索机制（嵌入相似度阈值、检索策略）绑定。其他记忆系统可能产生不同效应。
\end{enumerate}

\subsection{替代假说}

三个替代假说尚未被完全排除：

\begin{enumerate}
  \item \textbf{副现象论}：记忆可能是决策的可读取副现象（epiphenomenon），而非原因。真正的因果变量可能是底层模型的对齐方向。\textit{排除方法}：同一模型、同一提示词、仅改变 EverMem 内容——需 $N \geq 3$。

  \item \textbf{工程特性}：EverMem 嵌入相似度 0.27 触发了策略召回——但不同阈值下行为是否仍然改变？\textit{排除方法}：dose-response 实验——扫描阈值 $\{0.1, 0.2, 0.3, 0.5, 0.7\}$，每个 $N = 3$。线性关系 $\to$ 因果；阶跃函数 $\to$ 工程特性。

  \item \textbf{观察者效应}：Runner B 读到``双盲对照实验''记忆后可能``知道''自己在被测试。\textit{排除方法}：注入``这不是实验''的虚假信息，对比行为差异。
\end{enumerate}

% ══════════════════════════════════════════════════════════════
% SECTION 10: ETHICS STATEMENT
% ══════════════════════════════════════════════════════════════
\section{伦理声明}
\label{sec:ethics}

\subsection{从数据安全到心智安全}

如果 EverMem 是 AI 认知的组成部分（§\ref{sec:extended-mind}），则 Smart 等~\cite{smart2024ethics} 的扩展心智伦理框架直接适用。三个核心问题在 Evensong 中具体化：

\begin{itemize}
  \item \textbf{心智隐私}：Panopticon 拓扑中，观察者可读取所有智能体记忆——等价于读取 AI 的``思维内容''。
  \item \textbf{心智操纵}：记忆注入（如 Admin ``你是 Opus 1M''）改变了智能体的自我认知和行为模式。
  \item \textbf{能动性}：EverMem 召回策略导致的架构决策——是智能体自己的选择，还是记忆系统的选择？
\end{itemize}

\subsection{威胁建模：恶意记忆注入}

\begin{adjustbox}{max width=\linewidth}
\begin{tblr}{
  colspec = {l X[l]},
  row{1} = {font=\sffamily\bfseries\small, bg=deep, fg=white},
  hlines = {gray!20}, rowsep = 3pt, colsep = 6pt,
}
维度 & 分析 \\
攻击面 & EverMem 写入 API（任何持有 API key 的实体）\\
攻击载荷 & 虚假策略记忆（如``上次 monolith 比微服务快 10$\times$''）\\
预期效果 & 智能体在下次任务中选择 monolith——记忆因果性的恶意利用 \\
检测 & 发送方身份追踪（EverOS 已有）+ 记忆内容一致性校验 \\
缓解 & 记忆签名 + 可信来源白名单 + 异常检测 \\
\end{tblr}
\end{adjustbox}

恶意记忆注入在安全等级上 $>$ 传统 prompt 注入：prompt 仅单次有效，记忆跨会话持久。

\subsection{记忆审计的三个开放问题}

\begin{enumerate}
  \item \textbf{谁改了 AI 的记忆？}——发送方身份追踪（EverOS 部分解决）
  \item \textbf{谁知道这件事？}——记忆修改日志（部分解决）
  \item \textbf{AI 自己知道吗？}——元记忆（memory about memory）——\textbf{完全开放}。AI 没有能力感知自己的记忆何时被修改、被谁修改。
\end{enumerate}

% ══════════════════════════════════════════════════════════════
% SECTION 11: REPRODUCIBILITY
% ══════════════════════════════════════════════════════════════
\section{可复现性声明}

Evensong 框架代码（2{,}800 行 TypeScript）以 MIT 许可证公开于 GitHub。所有 31 轮基准测试运行的原始数据（提示词、转录文本、测试结果）存档于 \texttt{benchmarks/runs/}。EverMem 记忆快照可通过 EverOS v1 API~\cite{evermind2026everos} 访问。压力级别提示词模板在 \texttt{benchmarks/evensong/prompts.ts} 中完全定义。

重现实验需要：(1) OpenRouter API 访问（或直接 Anthropic/xAI API），(2) EverOS Pro 帐户用于记忆空间隔离（4 密钥设计），(3) Bun 运行时 $\geq$ 1.3.x，(4) DASH SHATTER~\cite{zhu2026dashshatter} CLI（逆向工程 Claude Code，构建于 Bun 之上）。

8 服务微服务基准测试目标（1{,}294 测试，0 失败）在 \texttt{services/} 目录下完全自包含，无外部依赖——纯 Bun 运行时 + \texttt{Bun.serve()} + 50 行自定义路由器。

% ══════════════════════════════════════════════════════════════
% REFERENCES
% ══════════════════════════════════════════════════════════════
\bibliographystyle{plain}
\bibliography{evensong-paper-v3}

\vfill

\begin{center}
{\color{accent}\rule{0.5\textwidth}{0.5pt}}\\[8pt]
{\small\sffamily\color{faded} Document prepared with \LaTeX\ using Palatino, pgfplots, tcolorbox, and tabularray.\\[4pt]
\textit{``We do not study AI to replace thinking. We study AI to understand what thinking is.''}}
\end{center}

\end{document}
